\SourceLeaf{82}%
\[
1+d+E+367(A-1700)-2b
\]
oder
\[
d+E+367(A-1700)-2b
\]
Ich ziehe ferner ab $364(A-1700)$, so kommt
\[
d+E+3A-5099-2b
\]
oder
\[
d+E+3A-5100-2b
\]
Ferner 5096 addirt gibt
\[
d+E+3A-3-2b
\]
oder
\[
d+E+3A-4-2b
\]
Endlich $3A-3c$, welches offenbar durch 7 theilbar ist, abgezogen gibt
\[
d+E+3c-3-2b
\]
oder
\[
d+E+3c-4-2b
\]
Dies von $7c+7d$ abgezogen, kommt
\[
3+2b+4c+6d-E
\]
oder
\[
4+2b+4c+6d-E
\]
welches also durch 7 theilbar sein muss. Hieraus ist klar, dass $E$ der Rest
sein werde, den man erhält, wenn man
\[
3+2b+4c+6d
\]
oder
\[
4+2b+4c+6d
\]
mit 7 dividirt, folglich $E=e$.

Es fällt also Ostern auf den $22+d+e^{\mathrm{ten}}$ März, oder (welches einerlei ist)
auf den $d+e-9^{\mathrm{ten}}$ April. W. Z. B. W.

\SourceLeaf{83}%
\begin{center}\textit{Ganz allgemeine Vorschriften zur Berechnung des Osterfestes
sowohl nach dem Julianischen, als nach dem Gregorianischen Kalender.}\end{center}
\begin{center}
\begin{tabular}{l|r|l}
Es entstehe aus der Division & mit & der Rest\\\hline
der Jahrzahl & 19 & $a$\\
der Jahrzahl & 4 & $b$\\
der Jahrzahl & 7 & $c$\\
der Zahl $19a+M$ & 30 & $d$\\
der Zahl $2b+4c+6d+N$ & 7 & $e$\\\hline
\end{tabular}\par
so fällt Ostern den $22+d+e^{\mathrm{ten}}$ März\par
oder den $d+e-9^{\mathrm{ten}}$ April
\end{center}

$M$ und $N$ sind Zahlen, die im Julianischen Kalender auf immer, im Gregorianischen
hingegen allemal wenigstens 100 Jahre hindurch unveränderliche
Werthe haben; und zwar ist in jenem $M=15$, $N=6$; in diesem, von der Einführung
desselben bis 1699, $M=22$, $N=2$
\begin{center}\small
\setlength{\tabcolsep}{5pt}\begin{tabular}{rrll|rrll}
von 1700 & bis 1799 & $M=23$, & $N=3$ & von 2100 & bis 2199 & $M=24$, & $N=6$\\
1800 & $\ldots$1899 & $M=23$, & $N=4$ & 2200 & $\ldots$2299 & $M=25$, & $N=0$\\
1900 & $\ldots$1999 & $M=24$, & $N=5$ & 2300 & $\ldots$2399 & $M=26$, & $N=1$\\
2000 & $\ldots$2099 & $M=24$, & $N=5$ & 2400 & $\ldots$2499 & $M=25$, & $N=1$
\end{tabular}
\end{center}

Allgemein findet man im Gregorianischen Kalender die Werthe von $M$
und $N$ für irgend ein gegebenes Jahrhundert von $100k$ bis $100k+99$ durch
folgende Regel:

Es gebe
\[
k\text{ mit }\left\{\begin{smallmatrix}3\\4\end{smallmatrix}\right\}
\text{ dividirt die (ganzen) Quotienten }
\left\{\begin{smallmatrix}p\\q\end{smallmatrix}\right\}
\]
wobei auf die Reste keine Rücksicht genommen wird.
Dann ist
\[
\left\{\begin{smallmatrix}M\\N\end{smallmatrix}\right\}
\text{ der \textit{Rest}, den man erhält, wenn man }
\left\{\begin{smallmatrix}15+k-p-q\\4+k-q\end{smallmatrix}\right\}
\text{ mit }\left\{\begin{smallmatrix}30\\7\end{smallmatrix}\right\}\text{ dividirt.}
\]
\textit{Beispiel.} Für die 100 Jahre von 4700 bis 4799 ist $k=47$, $p=15$, $q=11$;
also $15+k-p-q=36$; $4+k-q=40$; also $M=6$, $N=5$. So ist z.\,B.
für das Jahr 4763
\SourceLeaf{84}%
\begin{center}\small
\begin{tabular}{r|r|l}
$a=13$ & $19a+M=253$ & $e=3$\\
$b=3$ & $d=13$ & Ostern den $13+3-9$ d.\,i. den 7. April\\
$c=3$ & $2b+4c+6d+N=101$ & \quad nach dem Greg. Kalender
\end{tabular}\par\medskip
\begin{tabular}{r|l}
\multicolumn{2}{l}{Nach dem Julianischen hingegen}\\
$19a+M=262$ & $e=2$\\
$d=22$ & Ostern den $22+2-9$ d.\,i. den\\
$2b+4c+6d+N=156$ & \quad 15. April
\end{tabular}
\end{center}

Von obigen Regeln finden im \textit{Gregorianischen Kalender} einzig und allein
folgende zwei Ausnahmen Statt.

I. Gibt die Rechnung Ostern auf den 26. April, so wird dafür \textit{allemal} der
19. April genommen. (z.\,B. 1609. 1989).

Man sieht leicht, dass dieser Fall nur dann vorkommen kann, wo die Rechnung
$d=29$ und $e=6$ gibt; den Werth 29 kann $d$ nur dann erhalten, wenn
$11M+11$ mit 30 dividirt einen Rest gibt, der \textit{kleiner} als 19 ist; zu dem Ende
muss $M$ einen von folgenden 19 Werthen haben
\[
\begin{gathered}
0,\ 2,\ 3,\ 5,\ 6,\ 8,\ 10,\ 11,\ 13,\ 14,\ 16,\ 17,\ 19,\ 21,\ 22,\ 24,\ 25,\ 27,\ 29
\end{gathered}
\]
II. Gibt die Rechnung $d=28$, $e=6$, und kommt noch die Bedingung
hinzu, dass $11M+11$ mit 30 dividirt einen Rest gibt, der kleiner als 19 ist,
so fällt Ostern nicht, wie aus der Rechnung folgt, auf den 25., sondern auf den
18. April. --- Man überzeugt sich leicht, dass dieser Fall nur in denjenigen Jahrhunderten
eintreten könne, da $M$ einen von folgenden acht Werthen hat:
\[
2,\ 5,\ 10,\ 13,\ 16,\ 21,\ 24,\ 29
\]
\textit{Diese zwei Ausnahmen abgerechnet, sind obige Regeln völlig allgemein.}
\par\bigskip\centerline{\rule{.24\textwidth}{.3pt}}\medskip
[Handschriftliche Bemerkung:]

$p$ wird bestimmt als Quotient bei der Division von $8k+13$ mit 25.
\begin{center}\small
\begin{tabular}{lrl}
Von Ostern bis Michaelis sind im Durchschnitt $\ .\ .$ & $173\frac{52}{95}$ & Tage\\
und (Ostern ein, Michaelis ausgeschlossen) $\ .\ .\ .$ & $25\frac{21}{95}$ & Sonntage\\
& $148\frac{31}{95}$ & Wochentage\\
Von Michaelis bis Ostern hingegen $\ .\ .\ .\ .\ .\ .$ & $164\frac{773}{950}$ & Wochentage
\end{tabular}
\end{center}
\par\bigskip\centerline{\rule{.24\textwidth}{.3pt}}
